\section{Improving the SCD network with the DDS}
\label{sec:ddstcv}
The strict real-time requirements of a sophisticated control system such as the SCD, described in \Cref{sec:scd}, ask for an implementation that is designed and configured for the task at hand in any of its parts. Operating system kernels must be compiled with patches that enable fine-grained real-time task scheduling policies. The software must be written with frameworks that guarantee real-time execution, such as MARTe2. The hardware must be powerful enough to handle the computational load of the control algorithms, reserving to them all the system resources they require (\emph{e.g.}, CPU, RAM, PCI bandwidth) in a deterministic fashion, cutting out every other task including system ones. When sampling times get to the microsecond precision declared in \Cref{sec:scd}, every component of the system must be carefully considered.

In this scenario, things are a bit complicated by the fact that SCD is \emph{distributed} in the sense that has been described in \Cref{ch:intro}. It is not a single system but an interconnection of multiple computational nodes each executing a predefined set of tasks encoded in specific software subsystem, also interfacing with dedicated hardware components when necessary. Each node cannot perform its tasks alone: for example, control nodes 02 and 03 cannot actuate the coil power supplies properly without the real-time data provided by the plasma equilibrium reconstruction, which runs on nodes 05 and 06, which in turn need the data from the plasma diagnostics running on other nodes. By this description, it is clear how the communication infrastructure among the nodes must be carefully designed and implemented to guarantee the real-time performance of the whole system. It must employ a technology that makes it fast, deterministic, and reliable.

Currently (see \cite{galperti_overview_2024}), the SCD nodes are interconnected by a real-time fiber optic network that implements a \emph{reflective memory} (RFM). In essence, reflective memories are shared memory systems that allow multiple machines to share a data area. The implementation of the RFM usually consists of a high-speed connection medium, in this case an optic fiber connection, and dedicated network cards to install on the machines forming the network. These cards contain the actual data area, and their device drivers map that area within the main memory of the host machine. The RFM implementation is such that each local copy of the shared data area is kept coherent with the others, so that every write operation is propagated to all the other copies. The timings that these operations might take can be very tight but, most importantly, they are deterministic. This makes the RFM a suitable technology for this kind of application scenario.

The RFM technology has been commercially available since the 1980s and SCD adopts Abaco Systems network cards \cite{abaco_systems_pcie-5565rc_nodate}. As stated in \Cref{sec:scd}, these cards and their fiber optic interconnection allow SCD to maintain sub-microsecond latencies when transferring data among the nodes with a throughput up to \SI{170}{\mega\byte/\second}. This implementation has allowed SCD to perform well in the last decade but it is now showing some concerning issues, especially in terms of scalability and maintainability. First, as stated in \cite{abaco_systems_pcie-5565rc_nodate}, the Abaco cards are reaching their End Of Life (EOL) status with a limited production phase started in 2022. This, together with their high price (around 8.000€/card), is creating procurement concerns at SPC for such a critical component of the SCD. Secondly, the throughput of the RFM network is not enough to handle the increasing amount of data that the SCD nodes must exchange. New diagnostics are being evaluated and added each year, with some of them capable of generating considerable amounts of data, such as the MANTIS camera system on Node 10. The problem becomes clear when one notes that the data area is shared among \emph{all} the nodes, so they have to use that space to exchange data coming from every device and subsystem all at the same time. Despite some access policies that have been implemented to deal with this without increasing latencies (see, again, \cite{galperti_overview_2024}), a saturation point has been reached and new diagnostics such as bolometry cannot currently be used in real-time control because their data cannot be shared among the nodes.

To allow further upgrades to the SCD increasing TCV's experimental capabilities, as well as to avoid ending up relying on an EOL component, a replacement solutions for the RFM network has become necessary. Various criteria have been applied to its selection, including the cost-effectiveness of relying again on dedicated hardware components when common network hardware such as Ethernet has evolved so much in the meantime, guaranteeing bandwidths that were simply unthinkable at the time in which the reflective memory was invented. Ideally, a properly configured, reserved Ethernet network is enough if it driven by a software protocol that is able to guarantee the same determinism and low latencies of the RFM. The Data Distribution Service, introduced in \Cref{sec:dds}, is a candidate technology that can provide these features and has been selected as a replacement for the RFM network in the SCD.

This section describes the design of a software component for the SCD architecture that integrates the DDS middleware for communication purposes. Then, it presents some preliminary results of its application on SCD nodes, showing promising performance that motivated the decision to proceed with the integration of the DDS in the SCD. These results prove how the distributed technologies that constitute the base of the DUA framework presented in \Cref{ch:dua} can be applied successfully even to a different context such as the control system of a tokamak device.
